\section{总结与延伸}

本讲从“Transformer 只处理 tokens”出发，梳理了多模态模型的主要路线。CLIP/SigLIP 学习 image-text alignment；LLaVA 用 projector 把连续视觉 features 注入 LM；OneVision 和 Qwen-VL 系列改进 resolution、position、curriculum 与 cross-layer fusion；Chameleon 则把图像离散化，实现统一 autoregressive generation。

七个关键结论是：

\begin{enumerate}
\item 多模态 architecture 的第一约束是 token budget，而不是模型图的复杂程度。
\item CLIP-like representation 擅长 semantics，但不等于细粒度视觉理解。
\item Vision encoder + projector + LM 已成为 VLM 的稳定模板，主要创新常来自数据。
\item Dynamic resolution 与 MRoPE 用更多 compute 换取细节和时空结构。
\item Multi-image/video 的核心难点是 curriculum、loss balance 与 context allocation。
\item 连续 visual tokens 适合理解；离散 tokens 更容易统一生成，但会损失信息。
\item 真正的 omni model 必须同时解决 input understanding、native generation、training stability 和跨模态 evaluation。
\end{enumerate}

课程主线到这里形成闭环：从 tokenization 开始，我们最终又回到 tokenization，只是这次对象从文本扩展到了图像和视频。

\subsection{拓展阅读}

建议依次阅读 ViT、CLIP、SigLIP、LLaVA、LLaVA-OneVision、Qwen-VL/Qwen2-VL/Qwen3-VL、VQ-VAE 与 Chameleon，并重点比较它们的 token count、position encoding、data stages 和 generation interface。

\begin{knowledgebox}{复现实验：做一次 visual-token budget 验收}
准备四个固定 slice：自然图像问答、高分辨率文档 OCR、多图关系推理、短视频事件定位。对同一模型分别设置低/中/高 visual-token budget，逐项记录输入分辨率、tile/frame 数、最终 visual tokens、首 token latency、峰值显存和任务准确率。读结果时不要只找平均最优点：OCR 应检查最小可读字号，视频应检查采样遗漏，多图应检查图间指代，自然图像则检查是否为高预算支付了无收益的 attention cost。最后再改变 data mixture 或 freeze schedule 重跑一个 slice，验证性能变化来自 representation、训练数据还是 serving budget，而不是把所有收益都归因于模型版本名。
\end{knowledgebox}

\begin{importantbox}{最小交付门槛}
一次可信的多模态实验至少同时交付：每个 slice 的准确率与失败样本、visual/text token 数、prefill latency 与显存、resolution/frame policy、训练或评测数据版本，以及 generation 是否由原生 decoder 还是外部工具完成。若缺少其中任一项，就无法区分模型能力、数据覆盖和系统预算三种来源。
\end{importantbox}

\end{document}
